\documentclass[a4paper,11pt]{article}
\usepackage[T2A]{fontenc}
\usepackage[utf8]{inputenc}
\usepackage[russian]{babel}
\usepackage{amsmath,amssymb,mathtools}
\usepackage{geometry}
\geometry{left=22mm,right=22mm,top=18mm,bottom=18mm}
\usepackage{booktabs,tabularx,array}
\usepackage{xcolor,tikz}
\usepackage{hyperref}
\usepackage{microtype}
\usepackage{parskip}
\definecolor{accent}{HTML}{087C86}
\definecolor{ink}{HTML}{26333C}
\definecolor{errorred}{HTML}{A9363C}
\hypersetup{colorlinks=true,linkcolor=accent,pdfauthor={Лёвин Артём Александрович},pdftitle={Ревью домашней работы — Николь Саркисьянц — 02.11.2026}}
\setlength{\parindent}{0pt}
\setlength{\parskip}{3pt}
\renewcommand{\arraystretch}{1.23}
\newcommand{\Hdr}[1]{\par\vspace{2pt}\noindent\textbf{#1}\quad}
\newcommand{\E}{\par\smallskip\noindent\textcolor{errorred}{\textbf{Ошибка.}}\ }
\newcommand{\K}{\par\smallskip\noindent\textcolor{accent}{\textbf{Ключевая идея.}}\ }
\newcommand{\I}{\par\smallskip\noindent\textcolor{accent}{\textbf{Итог.}}\ }
\newcommand{\Task}[2]{\clearpage\phantomsection\label{task:#1}\section*{Задание №\,#1}\textcolor{accent}{#2}\par\vspace{7pt}}
\newcommand{\Footer}{\par\vfill\begin{center}\footnotesize\color{ink}Лёвин Артём Александрович эксклюзивно для Николь Саркисьянц\end{center}}
\begin{document}
\thispagestyle{empty}
\begin{center}
\vspace*{31mm}
{\large\color{accent} МАТЕМАТИКА}\par
\vspace{21mm}
{\Huge Ревью домашней работы}\par
\vspace{14mm}
{\Large\color{accent} Подготовка к ЕГЭ по профильной математике}\par
\vspace{22mm}
{\large Николь Саркисьянц}\par
\vspace{17mm}
{\large 2 ноября 2026 г.}\par
\vfill
{\small Лёвин Артём Александрович}\par
\smallskip
{\small эксклюзивно для Николь Саркисьянц}\par
\vspace{12mm}
\begin{tikzpicture}
\draw[accent, line width=.7pt] (0,0) .. controls (1,.5) and (2,-.5) .. (3,0)
 .. controls (4,.5) and (5,-.5) .. (6,0);
\end{tikzpicture}
\end{center}
\clearpage
\section*{Сводная таблица}
\noindent\small Нажмите на номер, чтобы перейти к подробному разбору. В таблицу включены только задания, для которых установлена ошибка или существенная незавершённость решения.\normalsize
\par\vspace{6pt}
\noindent\small
\begin{tabularx}{\linewidth}{@{}>{\centering\arraybackslash}p{10mm}p{26mm}X p{25mm} p{29mm}@{}}
\toprule
\textbf{№}&\textbf{Тема}&\textbf{Суть затруднения}&\textbf{Ваш ответ}&\textbf{Правильно}\\
\midrule
\hyperref[task:1]{1}&Геометрия&Найдена половина основания&14,4&28,8\\
\hyperref[task:2]{2}&Векторы&Вместо длины использовано перемножение координат&455&25\\
\hyperref[task:8]{8}&Тригонометрия&Потерян минус у косинуса&3&$-3$\\
\hyperref[task:9]{9}&Производная&Неверно подсчитаны участки возрастания&1&7\\
\hyperref[task:10]{10}&Движение&Неверно округлена доля времени&0,42&0,43\\
\hyperref[task:11]{11}&Движение&Решение начато, но не завершено&не указан&3,5 км/ч\\
\hyperref[task:12]{12}&Степенная функция&Нет вычисления отношения&не указан&27\\
\hyperref[task:14]{14}&Логарифмы&Некорректная проверка попадания на отрезок&нет вывода по б)&$(\log_5 34)^2$\\
\hyperref[task:16]{16}&Показательные&Ошибки при умножении на общий знаменатель; нет ответа&не указан&$1<x^2<3$\\
\hyperref[task:17]{17}&Производная&Пропущен $\cos\alpha$ при раскрытии скобок&не указан&204 м\\
\bottomrule
\end{tabularx}
\normalsize
\vspace{15pt}
\textcolor{ink}{\textbf{Примечание.}} Незавершённые и зачёркнутые без попытки задания не оценивались как ошибочно решённые. Задания, которые удалось проверить и которые решены правильно, в это ревью не включены.
\clearpage
\thispagestyle{empty}
\vspace*{.30\textheight}
\begin{center}{\Huge\color{accent} Разбор ошибок}\end{center}
\Task{1}{Геометрия равнобедренного треугольника}
\Hdr{Условие.} В треугольнике $ABC$ известно, что $AC=BC=15$, $\cos A=0{,}96$. Найдите $AB$.
\Hdr{Ваше решение.} Вы использовали $\cos A=\dfrac{x}{15}$ и нашли $x=15\cdot0{,}96=14{,}4$. В ответе записано $14{,}4$.
\E В качестве основания принята длина его половины.
\Hdr{Где именно.} Вычисление $15\cdot0{,}96=14{,}4$ верно. Ошибка возникает при переходе от полученного катета прямоугольного треугольника к ответу $AB$.
\Hdr{Почему.} В исходном равнобедренном треугольнике угол $A$ расположен между $AB$ и $AC$, но весь треугольник не является прямоугольным. Для применения косинуса нужно провести высоту $CD$ к основанию $AB$. В равнобедренном треугольнике она одновременно является медианой, поэтому $AD=DB$.
\K Косинус угла $A$ в прямоугольном треугольнике $ACD$ позволяет найти $AD$; затем основание удваивается.
\Hdr{Исправление.} После шага $AD=14{,}4$ запишите $AB=2AD$.
\Hdr{Полное правильное решение.} Проведём $CD\perp AB$. Тогда $AD=DB$. Из прямоугольного треугольника $ACD$:
\[
\cos A=\frac{AD}{AC},\qquad AD=15\cdot0{,}96=14{,}4.
\]
Следовательно,
\[AB=AD+DB=2AD=2\cdot14{,}4=\boxed{28{,}8}.\]
\Hdr{Проверка.} $AB=28{,}8<15+15=30$, условие существования треугольника выполняется. Половина основания равна $14{,}4$, что согласуется с косинусом.
\I $AB=28{,}8$.
\Hdr{Задача на закрепление.} В равнобедренном треугольнике $AC=BC=10$, $\cos A=0{,}8$. Найдите $AB$.
\Footer

\Task{2}{Координаты и длина вектора}
\Hdr{Условие.} Даны $\vec a=(6;-3)$ и $\vec b=(3;1)$. Найдите длину вектора $3\vec a-11\vec b$.
\Hdr{Ваше решение.} Правильно найдены $3\vec a=(18;-9)$ и $11\vec b=(33;11)$. Затем координаты этих векторов перемножены; получено число $455$.
\E При вычислении длины вектора вместо покоординатного вычитания выполнено действие, похожее на скалярное произведение. Кроме того, записанное вычисление $594-99=455$ арифметически неверно.
\Hdr{Где именно.} Последний корректный шаг: $3\vec a=(18;-9)$, $11\vec b=(33;11)$. Далее следует найти \emph{разность} этих векторов: $(18-33;\,-9-11)$.
\Hdr{Почему.} Вектор $3\vec a-11\vec b$ имеет две координаты. Длина вектора $(p;q)$ равна $\sqrt{p^2+q^2}$. Скалярное произведение $18\cdot33+(-9)\cdot11$ описывает другую операцию и к вопросу о длине суммы или разности векторов напрямую не относится.
\K Сначала выполните векторные операции отдельно с каждой координатой, затем примените теорему Пифагора.
\Hdr{Исправление.} $3\vec a-11\vec b=(-15;-20)$.
\Hdr{Полное правильное решение.}
\[
\begin{aligned}
3\vec a&=(18;-9),&11\vec b&=(33;11),\\
3\vec a-11\vec b&=(18-33;\,-9-11)=(-15;-20).
\end{aligned}
\]
По формуле длины:
\[
\bigl|3\vec a-11\vec b\bigr|=\sqrt{(-15)^2+(-20)^2}
=\sqrt{225+400}=\boxed{25}.
\]
\Hdr{Проверка.} $15^2+20^2=25^2$, следовательно, вычисление согласуется с теоремой Пифагора.
\I Длина искомого вектора равна $25$.
\Hdr{Задача на закрепление.} Даны $\vec a=(3;2)$ и $\vec b=(1;-1)$. Найдите длину вектора $2\vec a-4\vec b$.
\Footer

\Task{8}{Знаки тригонометрических функций}
\Hdr{Условие.} Вычислите
\[6\sqrt6\sin\frac{23\pi}{6}\cdot\cos\left(-\frac{11\pi}{4}\right):\tg\frac{8\pi}{3}.\]
\Hdr{Ваше решение.} Подставлены $\sin\frac{23\pi}{6}=-\frac12$, $\cos(-\frac{11\pi}{4})=\frac{\sqrt2}{2}$, $\tg\frac{8\pi}{3}=-\sqrt3$. Получен ответ $3$.
\E Косинусу угла $-\frac{11\pi}{4}$ присвоен положительный знак.
\Hdr{Где именно.} Значения синуса и тангенса определены верно. Первое неверное равенство: $\cos(-\frac{11\pi}{4})=\frac{\sqrt2}{2}$.
\Hdr{Почему.} Косинус --- чётная функция: $\cos(-t)=\cos t$. При этом
\[\frac{11\pi}{4}=2\pi+\frac{3\pi}{4},\]
а во второй четверти косинус отрицателен. Следовательно, $\cos(-\frac{11\pi}{4})=-\frac{\sqrt2}{2}$.
\K После приведения угла к стандартному промежутку обязательно установите четверть и знак функции.
\Hdr{Исправление.} Замените $+\frac{\sqrt2}{2}$ на $-\frac{\sqrt2}{2}$ и пересчитайте знак всего выражения.
\Hdr{Полное правильное решение.}
\[
\begin{aligned}
\sin\frac{23\pi}{6}&=\sin\left(4\pi-\frac\pi6\right)=-\frac12,\\
\cos\left(-\frac{11\pi}{4}\right)&=\cos\frac{3\pi}{4}=-\frac{\sqrt2}{2},\\
\tg\frac{8\pi}{3}&=\tg\frac{2\pi}{3}=-\sqrt3.
\end{aligned}
\]
Значит,
\[
6\sqrt6\cdot\left(-\frac12\right)\cdot\left(-\frac{\sqrt2}{2}\right):(-\sqrt3)
=\frac{3\sqrt{12}}{2}:(-\sqrt3)=\boxed{-3}.
\]
\Hdr{Проверка.} Числитель положителен (два отрицательных множителя), делитель отрицателен. Ответ обязан быть отрицательным.
\I $-3$.
\Hdr{Задача на закрепление.} Вычислите $4\sqrt2\sin\frac{7\pi}{6}\cdot\cos\frac{5\pi}{4}:\tg\frac{5\pi}{4}$.
\Footer

\Task{9}{Знак производной по графику функции}
\Hdr{Условие.} На графике $y=f(x)$ отмечены $x_1,\ldots,x_{13}$. Укажите количество отмеченных точек, в которых $f'(x)>0$.
\Hdr{Ваше решение.} В записи зафиксировано условие $f'(x)>0$, но в ответе указано $1$.
\E При подсчёте учитывался лишь один участок, где график идёт вверх; остальные отмеченные точки на возрастающих фрагментах пропущены.
\Hdr{Где именно.} Условие $f'(x)>0$ записано верно. Неверен переход от него к подсчёту: число таких точек на рисунке больше единицы.
\Hdr{Почему.} Положительная производная означает, что \emph{при движении по графику слева направо} функция локально возрастает. Значение самой функции, в том числе отрицательное, значения не имеет. В вершинах, где касательная горизонтальна, производная равна нулю и такие точки не подходят.
\K Последовательно проходите по всем вертикальным пунктирным линиям и отмечайте знак наклона графика в каждой точке.
\Hdr{Исправление.} Положительная производная наблюдается в точках $x_1,x_3,x_5,x_6,x_7,x_{10},x_{12}$.
\Hdr{Полное правильное решение.} Считываем направления наклона с исходного графика:
\[
\begin{array}{c|ccccccc}
\text{Точка} &x_1&x_3&x_5&x_6&x_7&x_{10}&x_{12}\\ \hline
\text{Поведение}&\nearrow&\nearrow&\nearrow&\nearrow&\nearrow&\nearrow&\nearrow
\end{array}
\]
Каждая стрелка $\nearrow$ обозначает возрастание графика в соответствующей точке. Всего таких точек $7$.
\Hdr{Проверка.} Точки на падающих участках и точки с горизонтальной касательной не добавляются к ответу. Отметка расположения выше/ниже оси $Ox$ роли не играет.
\I Подходящих точек $\boxed{7}$.
\Hdr{Задача на закрепление.} На графике отмечены шесть точек. В первых двух график локально возрастает, в третьей касательная горизонтальна, в четвёртой и пятой график убывает, в шестой вновь возрастает. В скольких отмеченных точках производная положительна?
\Footer

\Task{10}{Доля времени и округление}
\Hdr{Условие.} Скорость описывается формулой $v(t)=2\sqrt2\sin\dfrac{\pi t}{6}$ см/с. Какую долю первых $3{,}5$ секунды скорость не превышала $2$ см/с? Ответ округлите до сотых.
\Hdr{Ваше решение.} Найден подходящий промежуток длиной $1{,}5$ с и записано $\frac{1{,}5}{3{,}5}=0{,}42$.
\E Допущена арифметическая ошибка при переводе дроби в десятичную запись и округлении до сотых.
\Hdr{Где именно.} Граница $t=1{,}5$ с найдена верно. Первое неверное равенство: $\frac{1{,}5}{3{,}5}=0{,}42$.
\Hdr{Почему.} $\frac{1{,}5}{3{,}5}=\frac37=0{,}428571\ldots$; третья цифра после запятой равна $8$, поэтому при округлении до сотых получаем $0{,}43$.
\K После решения тригонометрического неравенства проверьте длину промежутка, разделите её на полную длительность и округлите только в самом конце.
\Hdr{Исправление.} Оставьте длительность $1{,}5$ с и вычислите долю без преждевременного округления.
\Hdr{Полное правильное решение.} На $0\le t\le3{,}5$ величина $\frac{\pi t}{6}$ меняется от $0$ до $\frac{7\pi}{12}$. Решаем:
\[
\begin{aligned}
2\sqrt2\sin\frac{\pi t}{6}&\le2,\\
\sin\frac{\pi t}{6}&\le\frac{\sqrt2}{2}.
\end{aligned}
\]
На рассматриваемом промежутке аргумента это выполнено при
\[
0\le\frac{\pi t}{6}\le\frac\pi4
\quad\Longleftrightarrow\quad 0\le t\le1{,}5.
\]
Следующий участок выполнения неравенства начинается после $t=4{,}5$ с и потому вне интервала наблюдения. Доля времени:
\[\frac{1{,}5}{3{,}5}=\frac37\approx\boxed{0{,}43}.\]
\Hdr{Проверка.} $0{,}42\cdot3{,}5=1{,}47$, что не равно $1{,}5$; $0{,}43\cdot3{,}5=1{,}505$ --- корректное округлённое приближение.
\I $0{,}43$.
\Hdr{Задача на закрепление.} Скорость $v(t)=2\sqrt2\sin\dfrac{\pi t}{8}$ см/с. Какую долю первых $4{,}5$ секунды скорость не превышала $2$ см/с? Округлите до сотых.
\Footer

\Task{11}{Относительное движение двух судов}
\Hdr{Условие.} Два судна длиной $150$ и $105$ м движутся параллельными курсами в одном направлении. Сначала второе позади первого, и от кормы первого до носа второго $170$ м. Через $15$ мин второе впереди, и от кормы второго до носа первого $450$ м. На сколько км/ч скорость второго больше?
\Hdr{Ваше решение.} Начата таблица «длина --- скорость --- время» для двух судов. Вычисление относительного перемещения и итоговый ответ отсутствуют.
\E Не завершено составление модели относительного движения: неизвестно, какое расстояние второе судно выиграло у первого за $15$ минут.
\Hdr{Где именно.} Корректная исходная мысль --- обозначить скорости двух судов. Дальше нужно вычислить изменение \emph{взаимного положения} с учётом двух корпусов и двух промежутков между судами.
\Hdr{Почему.} Длины судов нельзя игнорировать: исходные $170$ м и конечные $450$ м измерены между разными частями корпусов. От носа второго до его кормы есть $105$ м, а от носа первого до его кормы --- $150$ м.
\K При обгоне относительный путь равен начальному промежутку плюс длины обоих судов плюс конечный промежуток.
\Hdr{Исправление.} Запишите $S_{\text{отн}}=170+150+105+450$ и разделите на $15$ минут.
\Hdr{Полное правильное решение.} За время обгона второе судно относительно первого проходит:
\[
S_{\text{отн}}=170+150+105+450=875\text{ м}=0{,}875\text{ км}.
\]
Поскольку оба движутся в одном направлении, относительная скорость равна разности скоростей. Переведём $15$ минут в часы:
\[
t=\frac{15}{60}=0{,}25\text{ ч}.
\]
Следовательно,
\[
v_2-v_1=\frac{S_{\text{отн}}}{t}=\frac{0{,}875}{0{,}25}=\boxed{3{,}5\text{ км/ч}}.
\]
\Hdr{Проверка.} За четверть часа выигрыш при разности скоростей $3{,}5$ км/ч равен $0{,}875$ км, точно как найдено из условия.
\I Второе судно быстрее на $3{,}5$ км/ч.
\Hdr{Задача на закрепление.} Два судна длиной $120$ м и $80$ м движутся в одном направлении. Сначала второе отстаёт: расстояние от кормы первого до носа второго $100$ м. Через $10$ мин второе опережает первое: от кормы второго до носа первого $300$ м. На сколько км/ч второе быстрее?
\Footer

\Task{12}{График показательной функции}
\Hdr{Условие.} На рисунке дан график $f(x)=a^x$. Найдите $\dfrac{f(3)}{f(9)}$.
\Hdr{Ваше решение.} По графику отмечены точки $(-2;3)$ и приблизительно $(-3;5)$; выписаны соотношения $3=a^{-2}$ и $5=a^{-3}$. Итоговое отношение не найдено.
\E Решение остановлено после считывания точек. Дополнительную точку с приблизительным значением графика нельзя использовать как точное независимое условие.
\Hdr{Где именно.} Корректное равенство $a^{-2}=3$ уже достаточно для ответа. Дальше требуется применить свойство деления степеней.
\Hdr{Почему.} $\frac{a^m}{a^n}=a^{m-n}$ при $a>0$. Следовательно, отношение $f(3)/f(9)$ зависит от $a^{-6}$. Это третья степень уже известной величины $a^{-2}$. Приблизительное чтение высоты графика при $x=-3$ здесь избыточно.
\K Ищите в требуемом выражении степень того показателя, значение которого дано точно.
\Hdr{Исправление.} Из $a^{-2}=3$ получите $a^{-6}=(a^{-2})^3$.
\Hdr{Полное правильное решение.} По точке $(-2;3)$:
\[f(-2)=3\quad\Longrightarrow\quad a^{-2}=3.\]
По свойствам степеней:
\[
\frac{f(3)}{f(9)}=\frac{a^3}{a^9}=a^{3-9}=a^{-6}
=(a^{-2})^3=3^3=\boxed{27}.
\]
\Hdr{Проверка.} График убывает, значит $0<a<1$, и $f(3)>f(9)$; отношение $27>1$ согласуется с рисунком.
\I $27$.
\Hdr{Задача на закрепление.} На графике функции $f(x)=a^x$ отмечена точка $(-2;4)$. Найдите $\dfrac{f(2)}{f(8)}$.
\Footer

\Task{14}{Логарифмическое уравнение и отбор корней}
\Hdr{Условие.} Решите
\[
\log_{0{,}25}^{2}\bigl(5^{\sqrt{x}}-2\bigr)
-2\log_2\bigl(5^{\sqrt{x}}-2\bigr)+3{,}75=0,
\]
а затем найдите корни из отрезка $[3;7]$.
\Hdr{Ваше решение.} Корректно введена замена $t=\log_2(5^{\sqrt x}-2)$ и найдены $t=5$, $t=3$. Получены выражения $(\log_5 34)^2$ и $(\log_5 10)^2$. Затем принадлежность отрезку проверялась преобразованием квадрата логарифма в логарифм квадрата; окончательный ответ для части б) не записан.
\E Ошибка в способе сравнения корней с границами отрезка и отсутствие завершённого отбора. Нельзя считать $(\log_5 34)^2=\log_5(34^2)$.
\Hdr{Где именно.} Все шаги до записи корней сохраняются. Первая ненадёжная проверка начинается с перехода от неравенства для $(\log_5 34)^2$ к сравнению логарифмов от $125$, $1156$ и $78125$.
\Hdr{Почему.} Для любого положительного $b$:
\[
\log_5(b^2)=2\log_5 b,
\]
а квадрат $\bigl(\log_5 b\bigr)^2$ --- совсем другое выражение. При таком переходе изменяется проверяемое число.
\K Подстановка решает уравнение; принадлежность каждого найденного $x$ отрезку проверяется отдельно, с сохранением скобок у логарифмов.
\Hdr{Исправление.} Проверяйте $x_1=(\log_5 34)^2$ и $x_2=(\log_5 10)^2$ непосредственно.
\Hdr{Полное правильное решение.} Область определения:
\[5^{\sqrt x}-2>0,\quad x\ge0\quad\Longrightarrow\quad x>(\log_5 2)^2.\]
Поскольку $0{,}25=2^{-2}$, то $\log_{0{,}25}u=-\frac12\log_2u$. Пусть $t=\log_2(5^{\sqrt x}-2)$. Тогда
\[
\frac14t^2-2t+\frac{15}{4}=0
\quad\Longleftrightarrow\quad t^2-8t+15=0
\quad\Longleftrightarrow\quad(t-3)(t-5)=0.
\]
При $t=5$:
\[5^{\sqrt x}-2=32,\qquad \sqrt x=\log_5 34,\quad x_1=(\log_5 34)^2.\]
При $t=3$:
\[5^{\sqrt x}-2=8,\qquad \sqrt x=\log_5 10,\quad x_2=(\log_5 10)^2.\]
Проверим отрезок без округления. Так как $25<34<25\sqrt5$, имеем $2<\log_5 34<2{,}5$, следовательно $4<x_1<6{,}25<7$. Первый корень подходит. Далее $5<10<5\sqrt5$, откуда $1<\log_5 10<1{,}5$ и $1<x_2<2{,}25<3$. Второй корень не подходит.
\Hdr{Проверка.} Обе найденные величины положительны и удовлетворяют области определения; из них лишь $x_1$ лежит между $3$ и $7$.
\I В части а) два корня: $(\log_5 10)^2$ и $(\log_5 34)^2$. В части б) только $\boxed{(\log_5 34)^2}$.
\Hdr{Задача на закрепление.} Решите уравнение
\[
\log_{1/4}^{2}(2^{\sqrt{x}}-1)-2\log_2(2^{\sqrt{x}}-1)+3=0,
\]
после чего укажите корни, принадлежащие отрезку $[4;9]$.
\Footer

\Task{16}{Показательное неравенство: группировка}
\Hdr{Условие.} Решите
\[
2^{x^2-3}\bigl(3^{x^2-1}+4\bigr)<4^{x^2-2}+3^{x^2-1}.
\]
\Hdr{Ваше решение.} Вы раскрыли скобки, представили $4^{x^2-2}$ как степень двойки и перешли к дробям со степенями. На этапе избавления от знаменателей появились неверные коэффициенты; решение осталось без ответа.
\E Неправильное умножение каждого слагаемого на общий знаменатель и потеря удобной структуры исходного неравенства.
\Hdr{Где именно.} Верно:
\[
2^{x^2-3}\cdot3^{x^2-1}+4\cdot2^{x^2-3}-2^{2x^2-4}-3^{x^2-1}<0.
\]
При умножении дробной записи на $24$ коэффициент при $2^{x^2}$ должен стать $12$, а не $64$; при $2^{2x^2}$ --- $\frac32$, а не $8$.
\Hdr{Почему.} Если $2^{x^2-3}=2^{x^2}/8$, то
\[
24\cdot\frac{4\cdot2^{x^2}}8=12\cdot2^{x^2}.
\]
Кроме арифметики, раннее разложение всех степеней в дроби затрудняет поиск общего множителя.
\K Сгруппируйте отдельно слагаемые с основанием $3$ и с основанием $2$; обе группы содержат общий множитель $2^{x^2-3}-1$.
\Hdr{Исправление.} Обозначьте $A=2^{x^2-3}$. Тогда $4^{x^2-2}=4A^2$; это сохраняет все степени и резко упрощает запись.
\Hdr{Полное правильное решение.} Переносим всё в левую часть:
\[
\begin{aligned}
&2^{x^2-3}\cdot3^{x^2-1}+4\cdot2^{x^2-3}
-4^{x^2-2}-3^{x^2-1}<0.\end{aligned}
\]
Так как $4^{x^2-2}=4\bigl(2^{x^2-3}\bigr)^2$, получаем
\[
\begin{aligned}
&(2^{x^2-3}-1)3^{x^2-1}
 +4\cdot2^{x^2-3}(1-2^{x^2-3})<0,\\
&(2^{x^2-3}-1)\bigl(3^{x^2-1}-4\cdot2^{x^2-3}\bigr)<0,\\
&(2^{x^2-3}-1)\bigl(3^{x^2-1}-2^{x^2-1}\bigr)<0.
\end{aligned}
\]
Первый множитель отрицателен при $x^2<3$ и положителен при $x^2>3$. Второй имеет знак показателя $x^2-1$: при $x^2>1$ он положителен, при $x^2<1$ отрицателен (потому что $3^t$ и $2^t$ меняются местами при $t=0$). Произведение отрицательно только при
\[
1<x^2<3
\quad\Longleftrightarrow\quad
\boxed{x\in(-\sqrt3;-1)\cup(1;\sqrt3)}.
\]
\Hdr{Проверка.} При $x^2=2$ первый множитель отрицателен, второй положителен: подходит. При $x^2=0$ оба множителя отрицательны: не подходит. При $x^2=4$ оба положительны: не подходит. Границы исключаются, так как неравенство строгое.
\I $(-\sqrt3;-1)\cup(1;\sqrt3)$.
\Hdr{Задача на закрепление.} Решите неравенство
\[2^{x^2-2}\bigl(3^{x^2-1}+2\bigr)<2^{2x^2-3}+3^{x^2-1}.\]
\Footer

\Task{17}{Минимум функции: производная дроби}
\Hdr{Условие.} Минимизируйте расстояние
\[
l=h\frac{u+v\cos\alpha}{v\sin\alpha},\qquad h=85\text{ м},\ u=1{,}3\text{ м/с},\ v=0{,}5\text{ м/с},\quad 0<\alpha<\pi.
\]
\Hdr{Ваше решение.} После подстановки чисел получена функция
\[l(\alpha)=\frac{110{,}5+42{,}5\cos\alpha}{0{,}5\sin\alpha}.\]
Числитель и знаменатель продифференцированы правильно. При раскрытии числителя производной множитель $\cos\alpha$ при $110{,}5$ потерян. В результате получен ошибочный вывод о знаке производной; минимум и ответ не найдены.
\E Ошибка раскрытия скобок в формуле производной частного.
\Hdr{Где именно.} Правильно:
\[
U'= -42{,}5\sin\alpha,\qquad V'=0{,}5\cos\alpha.
\]
В $U'V-UV'$ член $-110{,}5\cdot0{,}5\cos\alpha$ нельзя заменить постоянным числом $-110{,}5$.
\Hdr{Почему.} Множитель $V'=0{,}5\cos\alpha$ умножает \emph{весь} числитель $U=110{,}5+42{,}5\cos\alpha$. Если отбросить косинус, меняется знак производной и исчезает точка минимума.
\K При дифференцировании частного сначала выпишите две производные и только затем раскрывайте $U'V-UV'$, не теряя множители.
\Hdr{Исправление.} Верный числитель производной:
\[
-21{,}25\sin^2\alpha-55{,}25\cos\alpha
-21{,}25\cos^2\alpha.
\]
\Hdr{Полное правильное решение.} По формуле производной частного:
\[
\begin{aligned}
l'(\alpha)
&=\frac{(-42{,}5\sin\alpha)(0{,}5\sin\alpha)
-(110{,}5+42{,}5\cos\alpha)(0{,}5\cos\alpha)}{0{,}25\sin^2\alpha}\\
&=\frac{-21{,}25(\sin^2\alpha+\cos^2\alpha)-55{,}25\cos\alpha}{0{,}25\sin^2\alpha}\\
&=-\frac{85+221\cos\alpha}{\sin^2\alpha}.
\end{aligned}
\]
Так как $\sin^2\alpha>0$ для $0<\alpha<\pi$, точка экстремума определяется равенством
\[
85+221\cos\alpha=0
\quad\Longleftrightarrow\quad
\cos\alpha=-\frac{85}{221}=-\frac5{13}.
\]
Тогда $\sin\alpha=\frac{12}{13}$. При меньших углах производная отрицательна, при больших положительна, значит найден минимум. Подставляем:
\[
\begin{aligned}
l_{\min}
&=85\cdot\frac{1{,}3+0{,}5\cdot(-5/13)}{0{,}5\cdot12/13}\\
&=85\cdot\frac{72/65}{6/13}
=85\cdot\frac{12}{5}=\boxed{204\text{ м}}.
\end{aligned}
\]
\Hdr{Проверка.} При приближении $\alpha$ к $0$ или к $\pi$ знаменатель стремится к нулю, а числитель остаётся положительным ($u>v$), поэтому расстояние растёт. Найденная внутренняя стационарная точка действительно даёт глобальный минимум.
\I Наименьшее возможное расстояние равно $204$ м.
\Hdr{Задача на закрепление.} Минимизируйте $l=h\dfrac{u+v\cos\alpha}{v\sin\alpha}$, если $h=45$ м, $u=1$ м/с, $v=0{,}6$ м/с и $0<\alpha<\pi$.
\Footer
\end{document}
